\cxtitle{Failure of the proposed Hamiltonian NEPv Rayleigh identity}

\begin{cxcredits}
\posedby{Amsel et al., \S6.2 \citeyearpar{AmselEtAl2026}}
\foundby{OpenAI Codex}{2026-08}
\formalizedby{OpenAI Codex}
\auditedby{OpenAI Codex}
\contributedby{Suvrit Sra}
\end{cxcredits}

\begin{cxcontext}
Section~6.2 considers $n,d\in\mathbb N$, nonzero vectors
$x_1,\ldots,x_d\in\mathbb C^n$, and matrices
$F_i,G_{i,j},K_{i,j}\in\mathbb C^{n\times n}$ that are Hermitian and have
norm at most one.  It defines
\[
H=\sum_{i=1}^{d}\left(F_i^{\langle i\rangle}
 +\sum_{j=1}^{d}G_{i,j}^{\langle i\rangle}
 K_{i,j}^{\langle j\rangle}\right)
\]
and the product-state Rayleigh quotient
\[
f(x_1,\ldots,x_d)=
\frac{(x_1\otimes\cdots\otimes x_d)^H
H(x_1\otimes\cdots\otimes x_d)}
{\prod_{j=1}^{d}x_j^Hx_j}.
\]
The proposed block diagonal matrix $A$ has blocks
\[
A_i=\left(\frac{\sum_{j=1}^{d}x_j^Hx_j}{x_i^Hx_i}\right)
\left(F_i+\sum_{j=1}^{d}G_{i,j}
\frac{x_j^HK_{i,j}x_j}{x_j^Hx_j}\right).
\]
The displayed source statement writes the stack through $x_n$, although $A$
has $d$ blocks and the surrounding text declares the NEPv dimension to be
$nd$.  The only dimensionally consistent reading is the stack through $x_d$;
the counterexample below uses $d=1$, so that $z=x_1$.
\end{cxcontext}

\begin{cxsource}{hamiltonian-nepv-identity}
The assertion appears immediately after equation~(20) in Section~6.2 of the
Simons workshop report \cite{AmselEtAl2026}.  That section is
scribed by Edgar Solomonik and records
discussion with Elias Jarlebring, Florian Schafer, Maryam Dehghan, Tamara
Kolda, and Chris Cama\~no; it says the question is based on a pre-workshop
question of Elias Jarlebring.
\end{cxsource}

\begin{cxstatement}{hamiltonian-nepv-identity}
Then, we have that with
\[
z=\begin{bmatrix}x_{1}\\
\vdots\\
x_{n}\end{bmatrix},\quad
\frac{z^{H}A(z)z}{z^{H}z}=f(x_{1},\ldots,x_{d}).
\]
\end{cxstatement}

\begin{cxsummary}{hamiltonian-nepv-identity}
The proposed NEPv matrix does not reproduce the product-state Rayleigh
quotient when the Hamiltonian sum includes a same-site term with $i=j$.
\end{cxsummary}

\begin{cxcertificate}{hamiltonian-nepv-identity}
The certificate consists of exact $2\times2$ integer matrices and the integer
vector $(1,2)^T$, for which the two assertedly equal values are $0$ and
$-4/25$.
\end{cxcertificate}

\begin{cxrefutation}
The construction in Section~6.2 of \cite{AmselEtAl2026} replaces a same-site
expectation of a product by the product of two expectations.  Those
quantities need not agree.

\begin{theorem}[\statusfalse: proposed NEPv Rayleigh identity]
\label{thm:hamiltonian-nepv-identity}
The displayed Rayleigh-quotient identity following equation~(20) is false for
the stated matrix formulas.  This refutes an assertion in the setup of
Problem~6.2, not the problem's request for an efficient algorithm.
\end{theorem}
\begin{proof}
Set $n=2$ and $d=1$, and take
\[
F_1=\begin{pmatrix}0&0\\0&0\end{pmatrix},\qquad
G_{1,1}=\begin{pmatrix}1&0\\0&0\end{pmatrix},\qquad
K_{1,1}=\begin{pmatrix}0&0\\0&-1\end{pmatrix},\qquad
x_1=\begin{pmatrix}1\\2\end{pmatrix}.
\]
All three matrices are Hermitian and have norm at most one.  Moreover,
$G_{1,1}$ and $K_{1,1}$ commute, so the matrix $H$ below is Hermitian.  Since
$d=1$, equation~(18) gives
\[
H=G_{1,1}K_{1,1}
=\begin{pmatrix}0&0\\0&0\end{pmatrix}.
\]
Therefore the objective in equation~(19) is
\[
f(x_1)=\frac{x_1^HHx_1}{x_1^Hx_1}=\frac{0}{5}=0.
\]
On the other hand, the scalar prefactor in equation~(20) is one and
\[
\frac{x_1^HK_{1,1}x_1}{x_1^Hx_1}=-\frac45,
\qquad
A(x_1)=-\frac45G_{1,1}.
\]
It follows that
\[
\frac{x_1^HA(x_1)x_1}{x_1^Hx_1}
=-\frac45\frac{1}{5}=-\frac{4}{25}\ne0=f(x_1).
\]
The failed step is specific and algebraic: for $i=j$, equation~(18) contains
\[
\frac{x_i^HG_{i,i}K_{i,i}x_i}{x_i^Hx_i},
\]
whereas equation~(20) produces
\[
\left(\frac{x_i^HG_{i,i}x_i}{x_i^Hx_i}\right)
\left(\frac{x_i^HK_{i,i}x_i}{x_i^Hx_i}\right).
\]
Hermiticity of the two factors does not make these expressions equal; the
counterexample uses commuting factors and even has $H=0$.
\end{proof}
\end{cxrefutation}
